% Ad Atticum 13.24 --- headnote
% ad-atticum-13-24, 11 July 45 BC, Tusculanum

\textit{Cicero to Atticus, written from the Tusculanum
on 11 July 709 AUC (Perseus dateline \textit{v Id.\
Quint.\ a.\ 709 (45)}), the day after Ad Atticum 13.23
and on the same two running threads: a rumoured
sighting of young Marcus Cicero at Corcyra on his way
to Athens, and the dedication of the
\textit{Academica} to Varro. Cicero has heard the
Corcyra report not from Atticus but at third hand,
through a freedman of Clodius reporting an Andromenes,
and is irritated that no letter has come from his son
even to his old friend. The four ``parchments''
\textit{[Greek: diphtherai]} are the four rolls of the
\textit{Academica} as now constituted, and they are in
Atticus' hands to release or withhold. Section 2 is a
brief follow-up to the previous day's letter on the
inheritance withholding-clause: Atticus is to push it
through without re-opening it.}

\textit{Two Greek tags carry the texture. The
\textit{[Greek: diphtherai]} is wonderfully concrete
--- four physical quires sitting on Atticus' table ---
and at the same time slightly distancing, as though
Cicero is trying to think of the dedication as a piece
of book-trade rather than the affectionate gesture he
has spent weeks worrying over. The
\textit{[Greek: aideomai]} (``I feel no scruple'')
disclaims, perhaps unconvincingly, the embarrassment
that has been driving the whole Varro thread; the
admission ``but I was the more afraid as to how he
himself would think of the thing'' is the real motive
just under the surface. The idiom \textit{in alteram
aurem} (``into the other ear with it'') is the Latin
counterpart of ``in one ear and out the other'' ---
Cicero is signalling that, having handed the
business off, he intends to stop fretting.}
